%% =====================================================================
%%  B5-T-14-08 -- what B5 contributes to "The Calibrated Question"
%%  -------------------------------------------------------------------
%%  LAYER A: APPEND-ONLY. Written once, then left alone. Material found
%%  only here sits in the `unique` slot and is protected by rule R10.
%% =====================================================================
%% @kind: source
%% @id: B5-T-14-08
%% @book: B5
%% @topic: T-14-08
%% @locator: ch. 7, pp. 188--215
%% @status: distilled
%% @unique: yes
%% @integrated: yes
%% @updated: 2026-09-19

\dossier{B5}{T-14-08}{\bookshort{B5}}{ch. 7, pp. 188--215}

\begin{distilled}
  \item The tool is the calibrated (open-ended) question: it \emph{removes aggression from conversations by acknowledging the other side openly, without resistance}, letting you introduce ideas and requests without sounding pushy --- a nudge.
  \item The allowed stems are \emph{how} and \emph{what}; \emph{why} is barred as an accusation in disguise; \emph{who/when/where} are fact-collectors, not steers.
  \item The polite-impossible refusal is a calibrated question: \emph{how am I supposed to do that?} --- sincerity of tone is the load-bearing part.
  \item The chapter names the mechanism the illusion of control: the counterpart does the solving, therefore owns the outcome, therefore defends it --- while the direction of the solution was set by the question.
  \item The Dos Palmas kidnapping case is the spine: calibrated questions used to hold a line with armed counterparts without escalating it.
  \item Standard inventory given for field use: what about this works for you; what is your objective; how can I help solve this problem; how would you like me to proceed.
\end{distilled}

\begin{unique}
  \item The grammar constraint as discipline: only how/what steer; everything else either collects facts or assigns blame --- which makes misfires diagnosable after the fact.
  \item The two-sided use stated openly: the same question that gives you control gives the counterpart the feeling of it, and the gap between the two is the entire technique.
  \item The refusal form as cooperation grammar: a no delivered as a how-question escalates nothing while conceding nothing.
\end{unique}
